\documentclass{article}
\usepackage{amsmath}
\title{Controller-Based Swing Optimization: Research Reference}
\begin{document}
\maketitle
\section{Scope and Conventions}
This editable research reference concerns the controller-based driver swing
optimization in issue 15493. Length, time, angle and torque are expressed in
metres, seconds, radians and newton metres unless otherwise stated. Coordinate
order, frames, model hashes and runtime must be recorded for each actual run.
The kinematic match is not itself a forward-dynamics validation.

\section{Control and Admission}
The intended controller is
\[
 \tau_{\rm cmd}=\tau_{\rm ff,ref}+\delta\tau_{\rm ff}
 +K_p(q_{\rm ref}-q)+K_d(\dot q_{\rm ref}-\dot q),
\]
with fixed gains and a smooth, bounded feedforward residual. Feasibility
must be checked against total applied torque, anatomical motion and speed,
including feedback and integration-stage behavior. Unavailable or assumed
limits must be identified; an incomplete envelope cannot certify feasibility.
The existing computed-torque tracker is a different controller.

\section{Contact and Flight}
Impact requires an approaching, bracketed signed-distance root on the finite
modeled face and the simulated event state $(t^*,q^*,\dot q^*)$. A face-centre
closest-approach heuristic is not that event. Missed contact must have a
finite smooth penalty. Prescribed-launch flight checks are separate from
swing-to-impact checks. Carry is distance to first ground contact and does not
include unvalidated rollout.
In the prototype, numerical trial and accepted substeps are partitioned on
the same fixed lattice anchored to each held-control interval, including an
arming split. This avoids a jump caused by changing the number of equal
substeps as the trial duration crosses a partition threshold. The bounded
scan can still miss narrow crossings; its candidate root is not a proof of
first contact, and the integration error is not enclosed.
On a folded mesh, a nearest triangle and its normal can change across a seam;
using that normal's side sign directly as a root function can create a false
discontinuous crossing. The prototype instead brackets finite-surface sphere
clearance, checks the outward-facing side separately, and requires a spatial
residual tolerance in addition to its time tolerance. This rejects the
demonstrated false root but does not close gaps in the bounded time scan.

\section{Evidence and Limitations}
The accompanying README records the prescribed-launch regression and the
separate tracking experiments. The renderer-pinned input embeds a calibrated
v1 model different from the committed v1 migration source. Its recorded
reference is a feedback rollout; the original tracking target is separate.
On the committed v2 model, the migrated initial state fails closure projection.
The largest initial residual component is 2.44519, in a mixed translation and
rotation residual vector; it is not a distance in metres. Matching coordinate
names does not establish model or closure equivalence. The exact model requires
an appropriately validated refit or another justified mapping. No accepted
optimized result is established. Extend this reference with geometry, solver
settings, sources, reproduction commands and margins as evidence is obtained.
\end{document}
